\chapter{Mise en place du pipeline Data Engineering}

\section{Introduction}
% TODO: Introduire le role du pipeline Data Engineering dans la solution.

\section{Ingestion des fichiers Excel}
% TODO: Decrire l'ingestion des fichiers Data Dump et Coverage.

\section{Traitement du Data Dump}
% TODO: Presenter les transformations principales liees aux visites, taches et reponses terrain.

\section{Traitement du Coverage}
% TODO: Presenter le role de fact_coverage pour la couverture, les statuts et les KPI.

\section{Modelisation et chargement dans MySQL}
% TODO: Decrire les tables principales et les relations utiles.

\begin{figure}[htbp]
    \centering
    \fbox{\parbox{0.82\textwidth}{\centering Espace reserve au diagramme du flux ETL.}}
    \caption{Flux general du pipeline ETL}
    \label{fig:flux_etl}
\end{figure}

\section{Tables dynamiques par tache}
% TODO: Expliquer le principe des tables task_* et leurs limites.

\section{Table normalisee survey\_responses}
% TODO: Expliquer pourquoi survey_responses facilite l'analyse et la validation.

\section{Integration future des Masters}
% TODO: Presenter les Masters comme extension future pour les dimensions de planification.

\section{Conclusion}
% TODO: Faire la transition vers la couche Business Intelligence.

